\chapter{向量组：相关性、基与坐标}\label{ch:vectors}
\chaptermap{相关与无关}{极大组与等价}{基与坐标}{内积与正交}
\begin{examtip}
\textbf{考研核心：}线性表示、相关性、极大无关组、等价向量组、基与坐标、Schmidt正交化。常见题把抽象向量组变成系数矩阵。
\textbf{竞赛视角：}用生成空间替代具体元素，用维数限制复杂关系。行变换保留列之间的线性关系，却通常不保留列空间的位置。
\end{examtip}
\section{线性组合与相关性的四种判定}
\subsection{从定义到矩阵方程}
\begin{definition}{生成、相关与无关}
设各向量属于同一$\F$上线性空间$V$。$\va$可由$\va_1,\ldots,\va_s$线性表示，是指存在$c_1,\ldots,c_s\in\F$使$\va=\sum_i c_i\va_i$；这些组合的全体为$\Span(\va_1,\ldots,\va_s)$。若关系$\sum_i c_i\va_i=0$必使全部$c_i=0$，称该组线性无关；若存在一组不全为零的系数使组合为零，称线性相关。含零向量的组必相关，单个非零向量无关。约定空向量组无关，其生成空间为$\{0\}$；全零向量组的极大无关子组可取空组，秩为0。
\end{definition}
\begin{theorem}{相关性判据与替换原则}
令$A=(\va_1,\ldots,\va_s)$，则列组无关当且仅当$\ker A=\{0\}$，当且仅当$\rank A=s$。当$s=n$且向量为$n$维时，又等价于$\det A\ne0$。无关组中加入$\vb$后相关，当且仅当$\vb$可由原组唯一表示。$n$维空间中超过$n$个向量必相关。
\end{theorem}
\begin{proof}
\label{der:3-1}
由列乘法定义，$A\vx=\sum_{j=1}^s x_j\va_j$。所以列组无关与$A\vx=0$只有零解是同一命题的两种写法。消元后若主元数$r<s$，令一个自由变量为1就产生非零解；若$r=s$，回代使全部未知量为零，故无关等价于满列秩。方形时满秩又等价于行列式非零，长方阵不使用这一判据。

若原组无关，加入$\vb$后有非零关系$\sum_i c_i\va_i+d\vb=0$。若$d=0$，原组无关迫使全部$c_i=0$，矛盾；故$d\ne0$，可解得$\vb=-\sum_i(\frac{c_i}{d})\va_i$。反向若$\vb=\sum_i b_i\va_i$，则$\sum_i b_i\va_i-\vb=0$是加入后的一条非零系数关系。若另有$\vb=\sum_i b_i'\va_i$，相减得$\sum_i(b_i-b_i')\va_i=0$，原组无关给$b_i=b_i'$，表示唯一。

最后，$n\times s$矩阵的主元最多$n$个；当$s>n$时不能满列秩，所以$n$维空间中超过$n$个向量必相关。

\end{proof}
\begin{method}{四法按维数选取}
定义法适合抽象向量；秩法适合长方阵；行列式法只用于向量数等于维数；齐次方程法可同时输出相关关系。要证明无关，需排除全部非零系数组合，不能只试几组。
\end{method}
\worked{一个相关关系的四种表达}{基础例题}{\star\star}{3-1}
{判定$\va_1=(1,0,1)\T,\va_2=(0,1,1)\T,\va_3=(1,1,2)\T$的相关性。}
{第三列逐项等于前两列之和。}
{\textbf{方法一（定义）：}$a_1+a_2-a_3=(1+0-1,0+1-1,1+1-2)\T=0$，系数$(1,1,-1)$非全零，所以相关。
\textbf{方法二（秩）：}前两列前两分量构成$I_2$，秩至少2；第三列等于前两列之和，秩至多2，所以秩2小于向量数3。
\textbf{方法三（行列式）：}对列矩阵作$c_3\leftarrow c_3-c_1-c_2$得到零列，行列式零，方形列组故相关。
\textbf{方法四（齐次方程）：}$c_1a_1+c_2a_2+c_3a_3=0$给$c_1+c_3=0,c_2+c_3=0,c_1+c_2+2c_3=0$。第三式冗余，取$c_3=t$得$(c_1,c_2,c_3)=t(-1,-1,1)$，$t=1$即非零关系。}
{四法描述同一结构；观察到显式关系时定义法最快。}
{“每两个都无关”不能推出三个向量无关。}
\worked{含参数的无关性}{考研风格}{\star\star}{3-2}
{判定$(1,0,0)\T,(1,t,0)\T,(1,1,t-1)\T$何时无关。}
{列矩阵为上三角。}
{\textbf{方法一：}拼列得到$A=\begin{pmatrix}1&1&1\\0&t&1\\0&0&t-1\end{pmatrix}$，上三角行列式为$t(t-1)$，所以$t\ne0,1$时无关，$t=0,1$时相关。
\textbf{方法二：}设关系系数$c_1,c_2,c_3$，逐分量为$c_1+c_2+c_3=0,tc_2+c_3=0,(t-1)c_3=0$。若$t\ne0,1$，从末式依次得$c_3=0,c_2=0,c_1=0$。若$t=0$，第一、二向量相同；若$t=1$，第二、三向量相同，均给显式非零关系。分类的必要性与充分性都已覆盖。}
{三角形结构使定义法和行列式法同样短。}
{不能只写充分条件而遗漏$t=0,1$的必要性。}
\worked{抽象向量的系数矩阵}{考研风格}{\star\star}{3-3}
{已知$\va,\vb$无关，判定$\va+\vb,\va-\vb$。}
{独立基向量可把关系转成系数关系。}
{设$x(a+b)+y(a-b)=0$，整理为$(x+y)a+(x-y)b=0$。原组$a,b$无关，故$x+y=0,x-y=0$，两式相加得$2x=0$，再得$x=y=0$，新组无关。另一条路线是系数矩阵$C=\begin{pmatrix}1&1\\1&-1\end{pmatrix}$，$\det C=-2\ne0$，因此可逆线性组合保持无关；两个论证都使用实数或复数域中$2\ne0$。}
{抽象向量不必设具体坐标，只需研究系数矩阵。}
{本题使用了$2\ne0$；这一条件在本册取用的实数域、复数域中成立。}
\worked{多项式向量组}{竞赛风格}{\star\star}{3-4}
{在次数不超过2的多项式空间中，证明$1,x+1,x^2+x$无关。}
{按最高次项逐层比较。}
{设$a\cdot1+b(x+1)+c(x^2+x)$恒等为零，合并同次项得到$cx^2+(b+c)x+(a+b)=0$。零多项式各系数必须为零，故$c=0,b+c=0,a+b=0$，依次得$a=b=c=0$。也可在单项式基$(1,x,x^2)$下写系数列矩阵$\begin{pmatrix}1&1&0\\0&1&1\\0&0&1\end{pmatrix}$，行列式1，故其列无关，与多项式向量组无关等价。}
{最高次项法是三角矩阵法的抽象版本。}
{“多项式等于零”表示恒等为零，不能只在一个点取值。}
\begin{insight}{相关是“有冗余”，不只是“平行”}
在平面上，两条不平行的向量已能生成全部方向；第三条向量即使不与前两条中的任何一条平行，仍是多余的。向量组相关的对象是整组系数关系，而不是只比较两两夹角。图中的$\va,\vb,\va+\vb$两两不平行，三个合起来却相关。
\end{insight}
\begin{figure}[H]\centering
\begin{tikzpicture}[scale=1.45]
\draw[laaxis](-.3,0)--(3.7,0)node[right]{$x$};\draw[laaxis](0,-.2)--(0,2)node[above]{$y$};
\draw[dashed,Teal!50](2,0)--(2.6,1.4)--(.6,1.4);
\draw[lavec,Ink](0,0)--(2,0)node[midway,below]{$\va$};
\draw[lavec,Gold](0,0)--(.6,1.4)node[above left]{$\vb$};
\draw[lavec,Teal](0,0)--(2.6,1.4)node[above right]{$\va+\vb$};
\node[lanote]at(2.3,-.65){$\va+\vb-(\va+\vb)=0$};
\end{tikzpicture}
\caption{第三条向量并未增加生成空间的维数。}\label{fig:dependent-vectors}
\end{figure}
\begin{clevertrick}{抽象向量题先“脱掉向量外衣”}
设$\va,\vb$无关，记新组$u=\va+\vb,v=\va-\vb,w=3\va+\vb$。任意关系$\alpha u+\beta v+\gamma w=0$整理成
\[(\alpha+\beta+3\gamma)\va+(\alpha-\beta+\gamma)\vb=0.\]
由于原组无关，两个标量系数必须同时为零。这正是二维系数列$(1,1)\T,(1,-1)\T,(3,1)\T$的相关性方程，所以可将抽象向量问题转成系数问题。前两列行列式为$-2\ne0$，故$u,v$无关。又$(3,1)\T=2(1,1)\T+(1,-1)\T$，从而
\[w=2u+v,\qquad 2u+v-w=0.\]
所以整组秩2，$u,v$是一个极大无关组，第三条的表示已完整给出，不必给$\va,\vb$设任何具体坐标。
\textbf{易错提示：}这一步保留无关性需要原向量组无关；若原组本来相关，系数矩阵满秩也不能保证新组无关。
\end{clevertrick}


\section{极大无关组与等价向量组}
\subsection{保留原列的身份}
\begin{definition}{极大无关组与等价}
原向量组中的一个无关子组，若加入原组任何其他向量都会相关，称极大无关组。其向量个数为该组的秩。两个向量组若互相可线性表示，称等价，即生成空间相同。
\end{definition}
\begin{theorem}{主元列与等价判定}
将列矩阵作初等行变换，阶梯形中的主元列编号对应原矩阵的一个极大无关组。若组$\mathcal A$可由组$\mathcal B$表示，则$\rank\mathcal A\le\rank\mathcal B$；在此基础上秩相等则两组等价。仅秩相等不足以等价。
\end{theorem}
\begin{proof}
\label{der:3-2}
设行阶梯形为$R=EA$，$E$可逆。对任何系数列$c$，$Rc=0\Longleftrightarrow Ac=0$，因此每一个列关系的系数在变换前后完全相同。阶梯形的主元列无关：从最后一个主元行向上考察，关系中的最后一个主元列系数先为零，再依次使所有系数为零。主元数为$r$，非主元列可通过从下向上的主元方程唯一表示为这$r$列的组合，因此它们生成全部列。把这些关系返回$A$，得到原主元列组成极大无关组。

若$\mathcal A$可由$\mathcal B$表示，则$U=\Span\mathcal A\subseteq W=\Span\mathcal B$。取$U$的基，它在$W$中仍无关，故$\dim U\le\dim W$。若维数相等而包含为真，取$w\in W\setminus U$，它能加到$U$的基中形成更长无关组，与同维矛盾；所以$U=W$，两组互相可表示。若只知维数相等却没有包含，两条不同直线就是反例。

\end{proof}
\begin{method}{求极大组并写全表示}
按列拼矩阵，作行阶梯形，记录主元列编号；回到原组取这些列；解出其他列的系数。易错：极大无关组一般不唯一，“极大”不表示挑选长度最大的向量。
\end{method}
\worked{找原向量组的基}{考研风格}{\star\star}{3-5}
{在例\ref{ex:3-1}的组中加入$\va_4=(2,1,3)\T$，求极大无关组及其余向量表示。}
{观察各列的前两个分量。}
{\textbf{方法一：}前两向量的前两个分量分别是$(1,0),(0,1)$，故独立。其余向量满足$a_3=a_1+a_2,a_4=2a_1+a_2$，所以它们生成全部四列，$\{a_1,a_2\}$为极大无关组。
\textbf{方法二：}列矩阵为$\begin{pmatrix}1&0&1&2\\0&1&1&1\\1&1&2&3\end{pmatrix}$，作$r_3\leftarrow r_3-r_1-r_2$得$\begin{pmatrix}1&0&1&2\\0&1&1&1\\0&0&0&0\end{pmatrix}$。主元编号1、2，第三列系数$(1,1)\T$、第四列$(2,1)\T$，返回原列得到同一表示。因为只有两个主元，秩恰为2。}
{观察法快，阶梯法能系统处理更多列。}
{答案应是原列$(1,0,1)\T,(0,1,1)\T$，不是消元后的单位列。}
\worked{同秩不等价}{基础例题}{\star}{3-6}
{判断组$\{(1,0)\T\}$与$\{(0,1)\T\}$是否等价。}
{比较生成直线。}
{第一组的全部组合为$\{(s,0)\T:s\in\R\}$，第二组为$\{(0,t)\T:t\in\R\}$。若第一组非零向量$(1,0)$能由第二组表示，会要求$(1,0)=(0,t)$，第一分量矛盾。反向也不能表示，因此两组不等价，虽然两组均只含一个非零向量，秩都为1。同秩只确认维数，不确认方向与集合相同。}
{秩只给维数，生成空间还包含方向信息。}
{若已知一组可由另一组表示，再结合秩相等才可推出等价。}
\worked{拼接矩阵判断表示}{考研风格}{\star\star}{3-7}
{设$A$的列组秩为$r$，给向量$\vb$，证明$\vb$可由$A$的列表示当且仅当$\rank(A,\vb)=r$。}
{增加一列时秩至多增加1。}
{\textbf{方法一：}若$b=Ac$，拼接列没有增加生成空间，故$\rank(A,b)=\rank A=r$。反向，原列空间包含于拼接列空间且二者维数相同，因此二者相等；新增列$b$必在原列空间，所以可表示。
\textbf{方法二：}求表示系数就是解$Ax=b$。行消元后，若增广列出现新的主元，会出现系数全零而右端非零的一行，方程无解且秩增加1；若不新增主元，两秩相等且回代可解。因此同秩条件正是可表示的充要条件。}
{空间法解释原理，方程法直接给系数。}
{不表示时秩必为$r+1$，但这仍不说明原列组是否独立。}
\worked{可逆系数组合保秩}{竞赛风格}{\star\star}{3-8}
{设$B=AP$且$P$可逆，证明$A,B$的列组等价。}
{写出反向表示。}
{逐列写$B=AP$：$b_j=\sum_i p_{ij}a_i$，所以$\Span(B)\subseteq\Span(A)$。$P$可逆使$A=BP\inv$，同样逐列得到反向包含，故两列组等价。若从像看，$\im(AP)=\{A(Px):x\in\F^s\}$；可逆$P$让$Px$遍历全部$\F^s$，因此该集合恰为$\im A$。这给同一结论的第二种证明，并指出奇异$P$时反向遍历不再成立。}
{可逆系数变换是“换生成元”，保留生成空间。}
{若$P$奇异，只能得到单向包含，未必等价。}
\section{基、坐标与过渡矩阵}
\subsection{固定一个方向约定}
\begin{definition}{有序基与坐标}
给定有限维线性空间$V$，既线性无关又生成$V$的有序组$\mathcal B=(\vb_1,\ldots,\vb_n)$称$V$的一组有序基。对每个$\vx\in V$，唯一表示$\vx=\sum_i x_i\vb_i$的系数列为$[\vx]_{\mathcal B}$。记$B=(\vb_1,\ldots,\vb_n),C=(\vc_1,\ldots,\vc_n)$，$C=BP$表示每列$\vc_j=\sum_i\vb_i p_{ij}$；对多项式等抽象向量，这是基列表的简记。约定$P$为从旧基$\mathcal B$到新基$\mathcal C$的过渡矩阵，其列为新基向量在旧基中的坐标。
\end{definition}
\begin{theorem}{坐标变换公式}
在$C=BP$约定下，
\begin{equation}[\vx]_{\mathcal B}=P[\vx]_{\mathcal C},\qquad[\vx]_{\mathcal C}=P\inv[\vx]_{\mathcal B}.\label{eq:coordinate}\end{equation}
有限维空间的任意两组基向量个数相同；该个数为维数。$n$维空间中的$n$个向量，只要无关或生成其一成立，就构成基。
\end{theorem}
\begin{proof}
\label{der:3-3}
记旧、新基矩阵为$B,C$且$C=BP$。同一向量分别写为$\vx=B\xi=C\eta=BP\eta$。因为旧基列无关，$B(\xi-P\eta)=0$推出$\xi=P\eta$。若$Pz=0$，则$Cz=BPz=0$，新基无关使$z=0$；$P$为方阵，故可逆，得到$\eta=P\inv\xi$。

为说明任意两组基个数相同，设长度为$m,n$。把第一组基用第二组表示，其系数列必须无关，否则对应的非零系数组合会使第一组相关；这些系数列处在$\F^n$，故$m\le n$。互换两组得$n\le m$，所以相等。若$n$维空间中有$n$个无关向量却不生成全空间，可加入一个不在其生成空间中的向量，得到$n+1$个无关向量，矛盾。若$n$个向量生成全空间却相关，删去一个可由其余表示的向量，剩下至多$n-1$个仍生成，会与空间维数$n$矛盾。因此两种单独条件都足以构成基。

\end{proof}
\begin{method}{基变坐标逆着变}
先写$C=BP$，不要只背“从哪到哪”；再从$B\xi=C\eta$推出$\xi=P\eta$。具体坐标可直接解$B\xi=\vx$，抽象基则比较系数。易错：过渡矩阵的列与坐标变换矩阵的方向不同。
\end{method}
\worked{非标准基下求坐标}{基础例题}{\star\star}{3-9}
{令$\mathcal B=((1,1)\T,(1,-1)\T)$，求$\vx=(3,1)\T$的坐标。}
{解两个系数的方程。}
{\textbf{方法一：}设$\vx=c_1(1,1)\T+c_2(1,-1)\T$，逐分量为$c_1+c_2=3,c_1-c_2=1$。相加得$c_1=2$，再得$c_2=1$，故坐标$(2,1)\T$。
\textbf{方法二：}$B=\begin{pmatrix}1&1\\1&-1\end{pmatrix}$，$\det B=-2$，$B\inv=\frac12\begin{pmatrix}1&1\\1&-1\end{pmatrix}$。乘原向量得$\frac12(3+1,3-1)\T=(2,1)\T$。代回$2(1,1)+(1,-1)=(3,1)$，表示核验完成。}
{一次求坐标可直接解方程，多向量时预求逆更方便。}
{向量本身$(3,1)\T$与它的坐标$(2,1)\T$不能混同。}
\worked{过渡矩阵与反向变换}{考研风格}{\star\star}{3-10}
{旧基为$\mathcal B=(\vb_1,\vb_2)$，新基为$(\vb_1+\vb_2,2\vb_2)$。若旧坐标为$(3,5)\T$，求新坐标。}
{把新基逐列写成旧基系数。}
{新基第一向量的旧坐标是$(1,1)\T$，第二向量是$(0,2)\T$，所以$C=BP$、$P=\begin{pmatrix}1&0\\1&2\end{pmatrix}$。设新坐标$\eta=(u,v)\T$，由旧坐标$P\eta=(3,5)\T$得$u=3,u+2v=5$，解得$v=1$。故新坐标$(3,1)\T$。直接展开$3(b_1+b_2)+1(2b_2)=3b_1+5b_2$，与旧表示相同。$\det P=2\ne0$也保证新组确为基。}
{代回展开可检验方向，比单纯记箭头更可靠。}
{新坐标不等于$P(3,5)\T$。}
\worked{连续换基的乘积}{考研风格}{\star\star}{3-11}
{若$C=BP,D=CQ$，求$\mathcal B$到$\mathcal D$的过渡矩阵及反向坐标公式。}
{保持矩阵乘法的顺序。}
{从$C=BP,D=CQ$依次代入，$D=(BP)Q=B(PQ)$，故从$B$到$D$的过渡矩阵是$PQ$。设旧坐标$\xi$、中间坐标$\eta$、最终坐标$\zeta$，同一向量给$\xi=P\eta,\eta=Q\zeta$，因此$\zeta=Q\inv P\inv\xi$。另一法直接将$\xi=PQ\zeta$左乘$(PQ)\inv$，而$(PQ)\inv=Q\inv P\inv$，得到同式。两个矩阵的先后顺序由基的定义决定。}
{基矩阵按实际定义连乘，坐标使用逆序逆矩阵。}
{$PQ$与$QP$通常不同。}
\worked{平移多项式基}{竞赛风格}{\star\star}{3-12}
{在$\mathcal P_2$中，以$\mathcal C=(1,x-1,(x-1)^2)$为基，求$x^2+1$的坐标。}
{令$y=x-1$，或比较单项式系数。}
{\textbf{方法一：}令$y=x-1$，则$x=y+1$，所以$x^2+1=(y+1)^2+1=y^2+2y+2$。按新基$(1,x-1,(x-1)^2)$顺序，坐标为$(2,2,1)\T$。
\textbf{方法二：}相对单项式基，过渡矩阵$P=\begin{pmatrix}1&-1&1\\0&1&-2\\0&0&1\end{pmatrix}$。解$P(u,v,w)\T=(1,0,1)\T$，从末行$w=1$，中行$v-2w=0$得$v=2$，首行$u-v+w=1$得$u=2$。三角矩阵对角全1可逆，所以坐标唯一。}
{变量代换最短，过渡矩阵法可同时处理多个多项式。}
{基是有序的，交换两个基向量时坐标也须交换。}
\section{内积、正交化与正交补}
\subsection{几何结构辅助计算}
\begin{definition}{内积与正交}
实线性空间上的内积是实值函数$\langle u,v\rangle$，满足对称性$\langle u,v\rangle=\langle v,u\rangle$、线性性$\langle\alpha u+\beta v,w\rangle=\alpha\langle u,w\rangle+\beta\langle v,w\rangle$和正定性$\langle u,u\rangle\ge0$，且等号仅在$u=0$时成立。由对称性可得另一变量也线性。

本章后续在$\R^n$的标准内积$\langle\vx,\vy\rangle=\vx\T\vy$下计算；对应欧氏范数为$\norm\vx_2=\sqrt{\langle\vx,\vx\rangle}$。内积为零称正交，长度为1称单位向量。单位正交组满足$\langle\vq_i,\vq_j\rangle=\delta_{ij}$。对指定内积空间$V$的子空间$W$，正交补是$W^\perp=\{\vx\in V:\langle\vx,\vw\rangle=0,\ \forall\vw\in W\}$，其含义依赖所在空间和内积。
\end{definition}
\begin{theorem}{Schmidt与Cauchy不等式}
无关组$\va_1,\ldots,\va_s$可按
\begin{equation}\vu_k=\va_k-\sum_{j<k}\dfrac{\langle\va_k,\vu_j\rangle}{\langle\vu_j,\vu_j\rangle}\vu_j,
\qquad\bm q_k=\dfrac{\vu_k}{\norm{\vu_k}_2}\label{eq:gs}\end{equation}
变为同生成空间的单位正交组。非零正交组必无关。另有$\abs{\langle\vx,\vy\rangle}\le\norm{\vx}_2\norm{\vy}_2$，等号当且仅当两向量相关。
\end{theorem}
\begin{proof}
\label{der:3-4}
假设已经得到非零且两两正交的$\vu_1,\ldots,\vu_{k-1}$。令$\vu_k=\va_k-\sum_{j<k}c_j\vu_j$，要求对每个$i<k$有$\langle\vu_k,\vu_i\rangle=0$，则
\[0=\langle\va_k,\vu_i\rangle-c_i\norm{\vu_i}^2,
\quad c_i=\dfrac{\langle\va_k,\vu_i\rangle}{\norm{\vu_i}^2}.\]
这推导出Schmidt公式。若$\vu_k=0$，则$\va_k$由先前$\vu_j$表示，而这些$\vu_j$由$\va_1,\ldots,\va_{k-1}$表示，违背原组无关，所以分母与单位化步骤均合法。公式正反解都为三角组合，故每一步新旧组生成空间相同。除以长度后，$\vq_i\T\vq_j=\delta_{ij}$。

非零正交组若有关系$\sum_i c_i\vu_i=0$，与$\vu_j$取内积得到$c_j\norm{\vu_j}^2=0$，故全部$c_j=0$，证明无关。

对Cauchy不等式，若$\vy=0$，两边均零。若$\vy\ne0$，将$\vx$减去沿$\vy$的投影：
\[\begin{aligned}
0&\le\left\|\vx-\dfrac{\vx\T\vy}{\vy\T\vy}\vy\right\|^2\\
&=\vx\T\vx-2\dfrac{(\vx\T\vy)^2}{\vy\T\vy}+\dfrac{(\vx\T\vy)^2}{\vy\T\vy}
=\norm\vx^2-\dfrac{(\vx\T\vy)^2}{\norm\vy^2}.
\end{aligned}\]
移项并乘正数$\norm\vy^2$、开非负平方根即得结论。等号当且仅当上述余量为零，即$\vx$为$\vy$的倍数；结合$\vy=0$情形，等号条件恰为两向量相关。

\end{proof}
\begin{method}{先正交后单位化}
逐步减去对已经得到的正交向量的投影；最后再除长度，通常减少根号。用点积为零检验每一步。若某步得到零，说明输入相关，应删去该向量；不能除以零。
\end{method}
\worked{Schmidt的完整两步}{考研风格}{\star\star}{3-13}
{正交化$\va_1=(1,1,0)\T,\va_2=(1,0,1)\T$。}
{第二个向量减去沿第一个方向的投影。}
{第一条$u_1=a_1=(1,1,0)\T$，$u_1\T u_1=2$。第二条的投影系数$\frac{a_2\T u_1}{u_1\T u_1}=\frac{1}{2}$，故
\[u_2=(1,0,1)\T-\dfrac12(1,1,0)\T=(\dfrac{1}{2},-\dfrac{1}{2},1)\T.\]
乘非零常数2，取同方向$u_2'=(1,-1,2)\T$。核验$u_1\T u_2'=1-1+0=0$，长度分别$\sqrt2,\sqrt6$，所以单位正交组为$q_1=\frac{(1,1,0)\T}{\sqrt2},q_2=\frac{(1,-1,2)\T}{\sqrt6}$。由$a_2=u_2+\frac{u_1}{2}$可反向表示，故生成空间未变。}
{正交化时允许非零倍乘，能避免分数。}
{第二步不能减去整个$\va_1$，系数必须由投影计算。}
\worked{正交组的无关性}{基础例题}{\star}{3-14}
{证明非零且两两正交的$\vu_1,\ldots,\vu_s$线性无关。}
{用内积逐个提取系数。}
{\textbf{方法一：}若$\sum_i c_iu_i=0$，对固定$j$取与$u_j$的内积，异指标项因正交消失，只剩$c_j\norm{u_j}^2=0$。$u_j\ne0$使平方长度正，所以$c_j=0$；对所有$j$成立，故无关。
\textbf{方法二：}令$U$以这些向量为列，则$U\T U=\diag(\norm{u_1}^2,\ldots,\norm{u_s}^2)$，对角元全正，矩阵可逆。若$Uc=0$，左乘$U\T$得$U\T Uc=0$，可逆性给$c=0$，再用定义证明列组无关。}
{内积法短，Gram法把几何转换成秩。}
{“非零”不能删除：零向量与所有向量正交，却使向量组相关。}
\worked{向直线投影}{考研风格}{\star\star}{3-15}
{求$\vb=(2,0)\T$在$\vu=(1,1)\T$生成直线上的正交投影。}
{投影为$c\vu$且残差与$\vu$正交。}
{\textbf{方法一：}投影点设$p=cu$。残差正交条件$u\T(b-cu)=0$给$c=\frac{u\T b}{u\T u}=\frac{2}{2}=1$，所以$p=(1,1)\T$，$e=b-p=(1,-1)\T$。核验$u\T e=1-1=0$。
\textbf{方法二：}任意候选点为$(c,c)$，距离平方$(2-c)^2+c^2=4-4c+2c^2=2(c-1)^2+2$。平方项非负，唯一最小在$c=1$，与正交法一致；最小距离为$\sqrt2$。}
{正交条件与距离最小化等价，第7章推广至子空间。}
{投影长度是$\sqrt2$，系数是$1$，两者不同。}
\worked{正交补的维数}{竞赛风格}{\star\star\star}{3-16}
{求$W=\{\vx\in\R^3:x_1+x_2+x_3=0\}$的正交补并证明维数相加为3。}
{方程的法向量是$(1,1,1)\T$。}
{\textbf{方法一：}由$x_3=-x_1-x_2$可选$W$的基$(1,0,-1)\T,(0,1,-1)\T$。若$y\perp W$，必须$y_1-y_3=0,y_2-y_3=0$，所以$y=t(1,1,1)\T$，反向代回对任意和零$x$内积为零。故$W^\perp=\Span((1,1,1)\T)$，维数$2+1=3$。
\textbf{方法二：}令$n=(1,1,1)\T$。任意$y$可写$y=\alpha n+z$，其中$\alpha=\frac{n\T y}{3}$且$z\in W$。若再有$y\in W^\perp$，则$0=y\T z=\alpha n\T z+z\T z=\norm z^2$，故$z=0$，所以$W^\perp=\Span(n)$。由约束消去$x_3$知$\dim W=2$，与$\dim W^\perp=1$相加为3；上述分解也直接给直和。}
{显式方程适合计算，四个子空间关系适合一般证明。}
{正交补包含零向量，不能只写一个单位法向量。}
\begin{insight}{Schmidt是在提取“新增的方向”}
把$\va_2$沿$\va_1$的部分减掉，剩余$\vu_2$才是第二个向量带来的新信息。若减完为零，表示$\va_2$没有新增方向；若非零，把它单位化，便得到一条与第一方向垂直的坐标轴。正交化改变基，却保留每一步已生成的空间。
\end{insight}
\begin{figure}[H]\centering
\begin{tikzpicture}[scale=1.5]
\draw[laaxis](-.4,0)--(3.1,0);\draw[laaxis](0,-.25)--(0,2.4);
\draw[lavec,Ink](0,0)--(2.5,0)node[right]{$\va_1$};
\draw[lavec,Gold](0,0)--(1.2,1.7)node[above]{$\va_2$};
\draw[lavec,Teal](0,0)--(0,1.7)node[above left]{$\vu_2$};
\draw[lavec,Teal!70](0,0)--(1.2,0)node[midway,below]{$p$};
\draw[dashed,Teal](1.2,0)--(1.2,1.7)--(0,1.7);
\draw[Ink!60](1.02,0)--(1.02,.18)--(1.2,.18);
\node[lanote]at(3.3,1.3){$p=\frac{\va_1\T\va_2}{\va_1\T\va_1}\va_1$\\[3pt]$\vu_2=\va_2-p$};
\end{tikzpicture}
\caption{投影部分$p$属于旧方向，余量$\vu_2$与旧方向正交。}\label{fig:gs-geometry}
\end{figure}
\begin{clevertrick}{先减投影，再放大去分母}
例\ref{ex:3-13}中$u_1=(1,1,0)\T$，第二条减投影得到$u_2=(\frac{1}{2},-\frac{1}{2},1)\T$。先取$w_2=2u_2=(1,-1,2)\T$，则
\[u_1\T w_2=1-1+0=0,\qquad\norm{w_2}^2=1+1+4=6.\]
故单位化结果$\frac{w_2}{\sqrt6}=\frac{(1,-1,2)\T}{\sqrt6}$。若直接单位化$u_2$，其长度为$\sqrt{\frac{1}{4}+\frac{1}{4}+1}=\frac{\sqrt6}{2}$，也得到同一向量。原方向与生成空间没有改变，因为$w_2=2u_2$与$u_2=\frac{w_2}{2}$可以互相表示。一般非零倍乘保持正交方向，能减少后续分数；若乘数为负，单位向量只反向，不影响单位正交基的合法性。
\textbf{易错提示：}单位化可以最后做，但下一次投影若使用尚未单位化的$\vu_j$，分母$\vu_j\T\vu_j$绝不能省略。
\end{clevertrick}


\begin{warning}
行消元不改变列之间的线性关系，但改变了列向量本身。最大无关子组有时有多种选择，秩却唯一。向量组的秩等于其生成空间的维数，未必等于整个所在空间$V$的维数。只有该组生成$V$时，它的极大无关子组才是$V$的一组基。正交方阵满足$Q\T Q=QQ\T=I$；矩形$Q$只有单位正交列时，$QQ\T$通常是投影而不是单位阵。
\end{warning}
\section*{章末分层习题}\addcontentsline{toc}{section}{章末分层习题}
\exercise{3-1}{基础}{判定$(1,2)\T,(2,4)\T$并写出非零相关关系。}
\exercise{3-2}{基础}{求$(1,0,0)\T,(0,1,0)\T,(1,1,0)\T$的极大无关组。}
\exercise{3-3}{基础}{求$(4,2)\T$在基$((1,1)\T,(1,-1)\T)$下的坐标。}
\exercise{3-4}{考研风格}{若$\va,\vb$无关，判定$\va+t\vb,2\va+\vb$何时无关。}
\exercise{3-5}{考研风格}{旧基$B$、新基$C=B\begin{pmatrix}2&1\\0&1\end{pmatrix}$，旧坐标$(5,1)\T$的新坐标是什么？}
\exercise{3-6}{考研风格}{将$(1,1,1)\T,(1,0,-1)\T$单位正交化。}
\exercise{3-7}{竞赛拓展}{证明无关组经过可逆线性变换仍无关。}
\exercise{3-8}{竞赛拓展}{在$\mathcal P_2$中证明$1,x,x^2$可由$1,1+x,1+x+x^2$唯一表示。}
\exercise{3-9}{竞赛拓展}{若$\norm{\va+\vb}^2=\norm{\va-\vb}^2$，证明$\va\perp\vb$。}
\section*{章末方法速查}\addcontentsline{toc}{section}{章末方法速查}
\begin{quicktable}
定义法&抽象向量&已知旧组无关&逐个系数比较\\
秩与齐次解&任意列数&元素具体&相关关系要非零\\
行列式法&方形列矩阵&向量数等于维数&长方阵不可用\\
主元列法&找极大无关组&消元后主元位置&回到原列\\
过渡矩阵&两组有序基&新基=旧基乘$P$&坐标用$P\inv$\\
Schmidt&无关输入&要求正交基&投影分母与零向量\\
\end{quicktable}
